\chapter{Pendahuluan}

\section{Latar Belakang}
Rapat, kegiatan lapangan, dan wawancara menghasilkan rekaman yang sering hanya disimpan, lalu jarang dibuka lagi. Ketika notulen atau laporan harus disusun, rekaman itu harus didengarkan ulang dari awal. Proses ini memakan waktu, dan kesalahan kecil dalam mencatat angka atau keputusan mudah terjadi. Di sisi lain, banyak instansi kini mencari cara agar hasil rekaman bisa dipakai sebagai bahan kerja yang cepat dan dapat diperiksa ulang.

Pandora dibangun untuk menjawab kebutuhan itu. Aplikasi ini mengubah rekaman audio menjadi teks, lalu membantu menyusun notulensi dan naskah dinas dari teks tersebut. Penggunaannya berada dalam kerangka tata kelola layanan digital pemerintah (Presiden Republik Indonesia, 2018). Penggunaan kecerdasan artifisial juga perlu mengikuti etika yang ditetapkan kementerian terkait (Kementerian Komunikasi dan Informatika Republik Indonesia, 2023).

\section{Deskripsi Singkat}
Modul ini membahas Pandora dari langkah pertama sampai hasil akhir. Bab II menjelaskan cara masuk dan mengenali bagian-bagian layar. Bab III membahas cara mentranskripsi rekaman. Bab IV membahas cara membaca hasil dan merapikannya. Bab V membahas notulensi rapat, dan Bab VI membahas dokumen resmi. Bab VII membahas riwayat hasil dan masalah yang sering muncul. Penutup mengajak Anda kembali ke pertanyaan pemantik di bab ini.

Setiap bab materi dibuka dengan indikator keberhasilan, dilanjutkan uraian yang disertai gambar atau tabel, lalu ditutup dengan latihan dan rangkuman. Untuk mengikat seluruh pembahasan, modul ini memakai satu ilustrasi kasus yang muncul di setiap bab. Ilustrasi itu bersifat rekaan, tetapi disusun dari jenis pekerjaan yang lazim dijumpai di lingkungan unit kerja.

\begin{ilustrasi}
Rina adalah staf sekretariat sebuah unit kerja. Pada suatu hari ia menghadiri tiga rapat, dan kepala unit meminta notulen ketiganya sudah siap sebelum rapat koordinasi besok pagi. Ia belum pernah memakai Pandora, dan ia tidak ingin salah mencatat keputusan yang dibaca banyak orang.

Bab-bab berikutnya mengikuti langkah Rina, mulai dari masuk ke aplikasi sampai ia menyusun notulen dan dokumen pertamanya.
\end{ilustrasi}

\section{Tujuan Pembelajaran}
\subsection{Hasil Belajar}
Setelah mengikuti pelatihan ini, peserta mampu menghasilkan transkrip, notulensi, dan dokumen resmi dari rekaman melalui Pandora, serta memeriksa hasil itu sebelum dibagikan.

\subsection{Indikator Hasil Belajar}
Peserta mampu:
\begin{enumerate}
  \item mengenali bagian layar Pandora dan membaca status layanan (Bab II);
  \item mentranskripsi satu rekaman dengan model dan bahasa yang sesuai (Bab III);
  \item membedakan transkrip asli dari transkrip rapi dan memeriksa angka terhadap transkrip asli (Bab IV);
  \item menyusun notulensi dan memeriksa keputusan, nama, serta tenggat terhadap transkrip (Bab V);
  \item menyusun dokumen resmi dengan kepala yang diisi sendiri dan memahami alasan naskah ditolak pemeriksa (Bab VI);
  \item mengelola riwayat hasil dan mengatasi masalah umum secara mandiri (Bab VII).
\end{enumerate}

\section{Materi Pokok dan Alokasi Waktu}
Modul ini dirancang untuk pelatihan selama 4 jam 20 menit dalam dua sesi. Tabel~\ref{tab:materi} memuat materi pokok setiap bab beserta perkiraan alokasi waktunya.

\begin{table}[H]
\caption{Materi pokok dan alokasi waktu}
\label{tab:materi}
{\small
\begin{tabularx}{\linewidth}{@{}K{1.6cm}YP{3.4cm}@{}}
\toprule
\rowcolor{ugmpanel}\thd{Bab} & \thd{Materi pokok} & \thd{Alokasi waktu} \\
\midrule
Bab I & Pendahuluan dan pertanyaan pemantik & 20 menit \\
Bab II & Masuk dan mengenal Pandora & 30 menit \\
Bab III & Transkripsi rekaman & 45 menit \\
Bab IV & Membaca dan merapikan hasil & 40 menit \\
Bab V & Notulensi rapat & 45 menit \\
Bab VI & Dokumen resmi & 45 menit \\
Bab VII & Riwayat dan masalah umum & 20 menit \\
Penutup & Kembali ke pertanyaan pemantik & 15 menit \\
\bottomrule
\end{tabularx}
}
\end{table}

\section{Petunjuk Belajar}
Setiap bab dibuka dengan indikator keberhasilan, lalu uraian materi yang diselingi ilustrasi kasus, latihan, dan rangkuman. Istilah teknis yang belum Anda kenal dapat dicari pada Glosarium di bagian awal modul. Untuk latihan, Anda memerlukan akun UGM yang sudah aktif di Pandora dan berkas contoh audio pendek yang disediakan pengelola pelatihan.

Informasi tentang fitur dalam modul ini diperiksa terhadap aplikasi pada 16 November 2026. Fitur dan pengaturan Pandora dapat berubah setelah tanggal itu. Bila tampilan di layar berbeda dengan modul, ikuti tampilan yang tersedia di aplikasi dan beri tahu pengelola agar modul diperbarui.

\section{Ketentuan Penggunaan Data}
Selama pelatihan, gunakan hanya dokumen publik atau rekaman contoh yang disediakan. Rekaman rapat internal, data pribadi, dan dokumen berklasifikasi tidak boleh diunggah sebelum ada persetujuan dari pimpinan unit dan kejelasan bahwa layanan yang dipakai sudah sesuai ketentuan. Ketentuan ini dibahas lagi pada Bab II dan Bab III, ketika Anda mengunggah rekaman pertama.

\section{Pertanyaan Pemantik}
Sebelum masuk ke materi, tuliskan jawaban Anda atas pertanyaan berikut, paling banyak lima butir. Jawaban yang sama akan Anda tinjau kembali pada Penutup.

\begin{quote}
\itshape ``Dari rapat atau kegiatan Anda minggu ini, apa yang paling sulit dicatat ulang? Siapa yang paling membutuhkan catatan itu? Bagaimana Anda tahu catatan itu benar?''
\end{quote}
